\documentclass[12pt, twoside]{article}
\input{../preamble}

\fancyhead[LO]{La Scuola d'Italia / Dr.\ Huson / IB Math 11 \\* 14 September 2026}
\fancyhead[RO]{\fbox{\begin{minipage}{6cm}\footnotesize First \& Last name:\\[0.5cm]\end{minipage}}}
\fancyfoot[C]{\thepage}

\begin{document}
\subsubsection*{1.2 Classwork: Arithmetic Sequences with Variables}

\begin{enumerate}[itemsep=0.15cm]

%--- Problem 1: Review -- negative d; u_n, u_20, solve u_n = -13 [7 marks] ---
%--- Source: Original (freshly authored, AA SL 1.2) ---
\item[1.]

Consider the arithmetic sequence
$$83,\; 77,\; 71,\; 65,\; \ldots$$

\begin{enumerate}[itemsep=0.1cm]
  \item Write down the value of the common difference $d$.
  \item Find an expression for $u_{n}$, the $n^{\text{th}}$ term.
  \item Find $u_{20}$.
  \item One term of the sequence is equal to $-13$. Find the value of
  $n$ for that term. Show your algebraic working.
\end{enumerate}

\begin{tikzpicture}
  \draw (-0.6,0) rectangle (14.6,14.7);
  \foreach \y in {1,2,...,14} \draw[dotted] (0,\y)--(14.1,\y);
\end{tikzpicture}

\newpage

%--- Problem 2: Terms in k; equal differences -> k, the terms, u_10 [7 marks] ---
%--- Source: Original (freshly authored, AA SL 1.2) ---
\item[2.]

The first three terms of an arithmetic sequence are
$$u_{1} = k - 4, \qquad u_{2} = 2k, \qquad u_{3} = 4k - 1.$$

\begin{enumerate}[itemsep=0.1cm]
  \item Find an expression, in terms of $k$, for $u_{2} - u_{1}$ and
  an expression for $u_{3} - u_{2}$. Simplify your answers.
  \item Use the fact that the sequence is arithmetic to find the value
  of $k$.
  \item Write down the first three terms of the sequence.
  \item Find $u_{10}$.
\end{enumerate}

\begin{tikzpicture}
  \draw (-0.6,0) rectangle (14.6,15.4);
  \foreach \y in {1,2,...,15} \draw[dotted] (0,\y)--(14.1,\y);
\end{tikzpicture}

\end{enumerate}
\end{document}
